% =============================================================
%  09_stile.tex
%  Capitolo 9 -- Stile e professione
% =============================================================

\chapter{Stile e professione}

\begin{intuizione}
Sai scrivere Lisp. Hai attraversato il modello mentale,
le funzioni, le liste, la ricorsione, gli higher-order,
le strutture dati, le macro, il backtracking, l'I/O.

\medskip
Questo capitolo è su qualcosa di diverso: come scrivere
\emph{buon} Lisp. Codice che si legge, si testa, si organizza
e si mantiene. Non una lista di regole, ma un modo di pensare
al codice come manufatto che sopravvive alla sessione di REPL.

\medskip
È anche il capitolo più soggettivo. Le scelte di stile
dipendono dal contesto. Quello che segue è un insieme
di principi con solide motivazioni, non dogmi.
\end{intuizione}

\vspace{4pt}

% ================================================================
\section{Funzioni piccole e API}
% ================================================================

La regola più importante del buon codice funzionale
non è specifica di Lisp: una funzione fa una cosa.

\begin{concetto}{Una funzione, una responsabilità}
\keyline{}{Una funzione dovrebbe essere descrivibile in una riga.}
\keyline{}{Se il nome ha una congiunzione (\emph{e}, \emph{o}, \emph{poi}),
probabilmente fa due cose.}
\keyline{}{Se il corpo è più lungo di una schermata, probabilmente
contiene più funzioni.}

\boxsep
\detail{Questa non è una regola sul numero di righe.
È una regola sulla coerenza concettuale.
Una funzione di 30 righe che fa davvero una cosa
è meglio di tre funzioni di 10 righe con responsabilità sovrapposte.}
\end{concetto}

\begin{esempio}[Refactoring: una funzione grande in tre]
\begin{lstlisting}
;; Versione monolitica: carica, filtra, formatta in una sola funzione.
(defun report-utenti-attivi (percorso)
  (with-open-file (stream percorso)
    (let ((risultato '()))
      (loop for riga = (read-line stream nil nil)
            while riga
            do (let* ((parti    (tokenizza riga))
                      (nome     (first parti))
                      (accessi  (parse-integer (second parti)))
                      (attivo-p (> accessi 10)))
                 (when attivo-p
                   (push (format nil "~a: ~a accessi" nome accessi)
                         risultato))))
      (nreverse risultato))))

;; Versione a tre funzioni: ogni livello ha una responsabilita'.
(defun carica-utenti (percorso)
  "Restituisce lista di (nome . accessi) dal file."
  (with-open-file (stream percorso)
    (loop for riga = (read-line stream nil nil)
          while riga
          collect (let ((parti (tokenizza riga)))
                    (cons (first parti)
                          (parse-integer (second parti)))))))

(defun utenti-attivi (utenti soglia)
  "Filtra gli utenti con accessi > soglia."
  (remove-if-not (lambda (u) (> (cdr u) soglia)) utenti))

(defun formatta-utente (utente)
  "Stringa leggibile per un singolo utente."
  (format nil "~a: ~a accessi" (car utente) (cdr utente)))

;; Composizione esplicita: si legge come prosa.
(defun report-utenti-attivi (percorso &optional (soglia 10))
  (mapcar #'formatta-utente
          (utenti-attivi (carica-utenti percorso) soglia)))
\end{lstlisting}
\end{esempio}

La versione a tre funzioni è più lunga, ma ogni parte
è testabile indipendentemente, riusabile separatamente
e comprensibile senza leggere il resto.

\begin{concetto}{API: il contratto di una funzione}
\keyline{Nome}{descrittivo, non abbreviato. \fn{conta-vocali} non \fn{cnt-v}.}
\keyline{Argomenti}{nell'ordine più naturale per l'uso parziale.
La ``cosa su cui si opera'' viene prima.}
\keyline{Valore di ritorno}{coerente: sempre dello stesso tipo,
o almeno sempre lo stesso tipo per lo stesso ramo.}
\keyline{Side effect}{dichiarati nel nome o nella docstring.
Se modifica lo stato globale, il nome dovrebbe suggerirlo.}
\end{concetto}

\begin{attenzione}
Non spezzare eccessivamente. Cinque funzioni di due righe
ciascuna, che si chiamano in sequenza senza possibilità
di riuso, sono peggio di una funzione di dieci righe.
Il punto è la \emph{separazione delle responsabilità},
non la minimizzazione delle righe.
\end{attenzione}

% ================================================================
\section{Nomi e leggibilità}
% ================================================================

In Common Lisp i nomi seguono alcune convenzioni forti.
Non sono obbligatorie, ma il codice che le rispetta
è immediatamente riconoscibile come ``idiomatico''.

\begin{concetto}{Convenzioni di naming}
\keyline{\fn{predicato-p}}{funzione che restituisce vero o falso:
\fn{null}, \fn{listp}, \fn{attivo-p}, \fn{valido-p}.}
\keyline{\fn{operazione!}}{funzione con side effect distruttivi:
\fn{sort}, \fn{nreverse}, \fn{delete}.
Il \fn{!} è una convenzione di stile, non sintassi.}
\keyline{\fn{*variabile*}}{variabile globale speciale (``ear muffs''):
\fn{*standard-output*}, \fn{*env-default*}.}
\keyline{\fn{+costante+}}{costante globale: \fn{+max-tentativi+}.}
\keyline{\fn{parola-parola}}{parole separate da trattino, non underscore
e non camelCase. \fn{conta-vocali}, non \fn{contaVocali}.}
\end{concetto}

\begin{confronto}
\textbf{Naming in Java vs Lisp.}
Java usa camelCase per metodi e variabili (\texttt{countVowels}),
PascalCase per classi (\texttt{UserAccount}).
Common Lisp usa kebab-case per tutto, con le convenzioni
\fn{-p}, \fn{!}, \fn{*...*}, \fn{+...+} per indicare il tipo
di binding.
Il reader di Common Lisp converte tutto in maiuscolo internamente,
quindi \fn{conta-vocali} e \fn{CONTA-VOCALI} sono lo stesso simbolo.
\end{confronto}

\argomento{Docstring}

\begin{concetto}{Docstring in \fnt{defun}}
\keyline{}{La prima stringa nel corpo di \fn{defun}, se presente,
è la docstring.}
\keyline{\fn{documentation}}{recupera la docstring a runtime:
\fn{(documentation 'nome-funzione 'function)}.}
\keyline{\fn{describe}}{mostra il tipo, la docstring e altri dettagli
di qualsiasi simbolo.}

\boxsep
\detail{Scrivi la docstring come se chi la legge conosce Lisp
ma non conosce la tua funzione. Prima frase: cosa fa.
Poi, se necessario: cosa si aspetta dagli argomenti
e cosa restituisce.}
\end{concetto}

\begin{esempio}[Docstring e \fn{describe}]
\begin{lstlisting}
(defun inverti-parole (stringa)
  "Restituisce una nuova stringa con le parole in ordine inverso.
   STRINGA deve essere una stringa con parole separate da spazi.
   Restituisce una stringa; se STRINGA e' vuota restituisce \"\"."
  (let ((parole (tokenizza stringa)))
    (format nil "~{~a~^ ~}" (nreverse parole))))

;; Recupero della docstring a runtime
(documentation 'inverti-parole 'function)
; => "Restituisce una nuova stringa con le parole in ordine inverso. ..."

;; describe mostra molte piu' informazioni
(describe 'inverti-parole)
; INVERTI-PAROLE names a compiled function
;   required args: (STRINGA)
;   ...docstring...
\end{lstlisting}
\end{esempio}

% ================================================================
\section{Package e namespace}
% ================================================================

Fino a qui tutto il codice è vissuto nel package
di default, \texttt{CL-USER}.
Appena il progetto ha più di un file, serve un sistema
di namespace per evitare conflitti di nomi.

\begin{concetto}{Package}
\keyline{}{Un package è uno spazio di nomi: un insieme di simboli
con un nome comune.}
\keyline{\fn{defpackage}}{definisce un package con i suoi import ed export.}
\keyline{\fn{in-package}}{cambia il package corrente.}
\keyline{\fn{export}}{rende un simbolo visibile a chi usa il package.}
\keyline{\fn{use-package}}{importa i simboli esportati da un altro package.}

\boxsep
\detail{I simboli non esportati sono privati: visibili solo
dall'interno del package. È il meccanismo di incapsulamento
di Common Lisp, equivalente al \texttt{private} di Java.}
\end{concetto}

\begin{esempio}[Due package che interagiscono]
\begin{lstlisting}
;;; File: matematica.lisp

(defpackage :matematica
  (:use :cl)
  (:export :fattoriale :fibonacci))

(in-package :matematica)

(defun fattoriale (n)
  "n! ricorsivo."
  (if (<= n 1) 1 (* n (fattoriale (1- n)))))

(defun fibonacci (n)
  "n-esimo numero di Fibonacci."
  (cond ((= n 0) 0) ((= n 1) 1)
        (t (+ (fibonacci (- n 1)) (fibonacci (- n 2))))))

;; helper interno (non esportato: non visibile dall'esterno)
(defun %valida-n (n)
  (check-type n (integer 0 *) "intero non negativo"))
\end{lstlisting}

\begin{lstlisting}
;;; File: main.lisp

(defpackage :mio-programma
  (:use :cl :matematica))  ; importa i simboli esportati da :matematica

(in-package :mio-programma)

(fattoriale 10)    ; => 3628800  (simbolo importato da :matematica)
(fibonacci 10)     ; => 55

;; Accesso qualificato (senza use-package):
(matematica:fattoriale 5)  ; => 120
\end{lstlisting}
\end{esempio}

\begin{attenzione}
Se due package esportano lo stesso simbolo e li importi
entrambi con \fn{use-package}, ottieni un conflitto.
La soluzione è accedere al simbolo con il nome qualificato
(\fn{package:simbolo}) invece di importarlo.
\end{attenzione}

% ================================================================
\section{Progetti multi-file e ASDF}
% ================================================================

Un progetto reale ha più file da compilare nell'ordine giusto.
ASDF (Another System Definition Facility) è lo strumento
standard per descrivere questa struttura.

\begin{concetto}{ASDF: cos'è e cosa fa}
\keyline{}{ASDF descrive un sistema: quali file lo compongono,
in che ordine compilarli, quali dipendenze esterne richiede.}
\keyline{System definition}{un file \fn{.asd} nella radice del progetto.}
\keyline{\fn{asdf:load-system}}{carica e compila il sistema nella
sequenza corretta.}
\keyline{Quicklisp}{il gestore di pacchetti che usa ASDF internamente.
\fn{(ql:quickload :nome-libreria)} scarica e carica una libreria.}
\end{concetto}

\begin{esempio}[Un file .asd minimale]
\begin{lstlisting}
;;; File: mio-progetto.asd

(asdf:defsystem :mio-progetto
  :description "Un progetto di esempio"
  :version "0.1.0"
  :author "Carmelo Calzerano"
  :license "MIT"
  :depends-on ()           ; librerie esterne necessarie
  :components
  ((:file "packages")      ; deve venire primo: definisce i package
   (:file "matematica" :depends-on ("packages"))
   (:file "main"       :depends-on ("packages" "matematica"))))
\end{lstlisting}
\end{esempio}

\begin{nota}
Questa sezione mostra la forma base.
ASDF ha molte altre opzioni: test runner integrato, versioni
dei file, gerarchie di moduli, sistemi condizionali.
La documentazione completa è su \url{https://asdf.common-lisp.dev}.
Per la maggior parte dei progetti piccoli e medi,
la struttura mostrata sopra è sufficiente.
\end{nota}

% ================================================================
\section{Testing}
% ================================================================

Scrivere test prima di dichiarare finita una funzione
non è solo buona pratica: è il modo più rapido
per capire se la funzione fa quello che pensi.

\begin{concetto}{Testing in Common Lisp}
\keyline{\fn{assert}}{macro della libreria standard: segnala un errore
se la condizione è falsa.}
\keyline{Framework}{librerie come \fn{fiveam} o \fn{parachute}
aggiungono suite di test, report, fixture.}
\keyline{Uso minimale}{per progetti piccoli \fn{assert} basta.
Per progetti che crescono, un framework evita di
reinventare la ruota.}
\end{concetto}

\begin{esempio}[Test con \fn{assert} e con \fn{fiveam}]
\begin{lstlisting}
;;; Test minimali con assert (nessuna dipendenza)
(defun testa-fattoriale ()
  (assert (= (fattoriale 0) 1))
  (assert (= (fattoriale 1) 1))
  (assert (= (fattoriale 5) 120))
  (assert (= (fattoriale 10) 3628800))
  (format t "fattoriale: tutti i test passati~%"))

(testa-fattoriale)
; fattoriale: tutti i test passati
\end{lstlisting}

\begin{lstlisting}
;;; Test con fiveam (da installare con Quicklisp)
;; (ql:quickload :fiveam)

(defpackage :mio-progetto/test
  (:use :cl :fiveam :mio-progetto))

(in-package :mio-progetto/test)

(def-suite suite-matematica
  :description "Test per le funzioni matematiche")

(in-suite suite-matematica)

(test fattoriale-casi-base
  (is (= (fattoriale 0) 1))
  (is (= (fattoriale 1) 1)))

(test fattoriale-valori-medi
  (is (= (fattoriale 5) 120))
  (is (= (fattoriale 10) 3628800)))

;; Esegui la suite
(run! 'suite-matematica)
; Running test suite SUITE-MATEMATICA
;  Running test FATTORIALE-CASI-BASE ..
;  Running test FATTORIALE-VALORI-MEDI ..
; Did 4 checks.
;    Pass: 4 (100%)
\end{lstlisting}
\end{esempio}

\begin{intuizione}
Scrivere il test prima di scrivere la funzione non è
masochismo: è il modo più veloce per chiarire
esattamente cosa la funzione deve fare.
Se non riesci a scrivere il test, non hai ancora capito
abbastanza bene il problema.
\end{intuizione}

% ================================================================
\section{Condizioni e gestione degli errori}
% ================================================================

Common Lisp ha un sistema di gestione delle condizioni
più potente del semplice try/catch.
Qui ne vediamo la forma base, che copre la maggior parte
dei casi pratici.

\begin{concetto}{Condizioni}
\keyline{\fn{error}}{segnala una condizione di errore non recuperabile
(salvo handler).}
\keyline{\fn{warn}}{segnala un avviso: il programma continua.}
\keyline{\fn{handler-case}}{intercetta condizioni, come try/catch.}
\keyline{\fn{ignore-errors}}{esegue il corpo; se c'è un errore
restituisce \fn{nil} e il messaggio d'errore.}
\keyline{\fn{check-type}}{segnala un errore di tipo se il valore
non è del tipo atteso.}
\end{concetto}

\begin{esempio}[Segnalare e intercettare condizioni]
\begin{lstlisting}
;; Funzione che segnala un errore al posto di crashare silenziosamente
(defun dividi (a b)
  "Divide A per B. Segnala un errore se B e' zero."
  (when (zerop b)
    (error "Divisione per zero: ~a / ~a" a b))
  (/ a b))

;; handler-case intercetta l'errore
(handler-case
    (dividi 10 0)
  (error (e)
    (format t "Errore catturato: ~a~%" e)
    nil))
; Errore catturato: Divisione per zero: 10 / 0
; => NIL

;; ignore-errors: forma breve per "prova, se fallisce restituisci nil"
(ignore-errors (dividi 10 0))   ; => NIL, "Divisione per zero: 10 / 0"
(ignore-errors (dividi 10 2))   ; => 5
\end{lstlisting}
\end{esempio}

\begin{confronto}
\textbf{Confronto con Java.}
In Java si usa \texttt{try}/\texttt{catch}/\texttt{finally}.
\fn{handler-case} è l'equivalente diretto.
La differenza principale è che Common Lisp distingue
tra condizioni segnalate (\fn{error}, \fn{warn}, \fn{signal})
e handler che le intercettano (\fn{handler-case},
\fn{handler-bind}).
Il sistema di condizioni è più flessibile del modello
exception-based perché permette di riprendere l'esecuzione
dal punto in cui è stata segnalata la condizione
(restart), senza necessariamente salire nello stack.
\end{confronto}

\begin{attenzione}
\fn{ignore-errors} è comodo ma pericoloso se usato
senza attenzione: nasconde tutti gli errori, anche
quelli che indicano bug nel codice.
Usarlo solo quando si sa esattamente quale condizione
può accadere e si è sicuri di volerla ignorare.
Per tutto il resto, preferire \fn{handler-case} con
una clausola esplicita.
\end{attenzione}

% ================================================================
\section{Profiling e ottimizzazione}
% ================================================================

Prima di ottimizzare, misura.
Il profiling è lo strumento che trasforma ``questo codice
sembra lento'' in ``questa funzione specifica consuma il 90\%
del tempo''.

\begin{concetto}{Strumenti di profiling in Common Lisp}
\keyline{\fn{time}}{macro standard: esegue la forma e stampa
il tempo trascorso e l'allocazione di memoria.}
\keyline{\fn{room}}{mostra l'uso corrente della memoria heap.}
\keyline{\fn{declare optimize}}{direttiva al compilatore:
più velocità, meno controlli di tipo.}

\boxsep
\detail{Per profiling più dettagliato (quale funzione specificamente),
Allegro CL offre \fn{prof:with-profiling} / \fn{prof:show-flat-profile}.
SBCL offre \fn{sb-profile:profile}. La sintassi varia per
implementazione.}
\end{concetto}

\begin{esempio}[\fn{time}: misura di una funzione ricorsiva]
\begin{lstlisting}
;; Fibonacci ingenuo: complessita' esponenziale
(defun fib-lento (n)
  (cond ((= n 0) 0) ((= n 1) 1)
        (t (+ (fib-lento (- n 1)) (fib-lento (- n 2))))))

(time (fib-lento 35))
; Evaluation took:
;   1.234 seconds of real time
;   ...
; => 9227465

;; Fibonacci con memoization: complessita' lineare
(let ((cache (make-hash-table)))
  (defun fib-veloce (n)
    (or (gethash n cache)
        (setf (gethash n cache)
              (cond ((= n 0) 0) ((= n 1) 1)
                    (t (+ (fib-veloce (- n 1))
                          (fib-veloce (- n 2)))))))))

(time (fib-veloce 35))
; Evaluation took:
;   0.000 seconds of real time
;   ...
; => 9227465
\end{lstlisting}
\end{esempio}

\begin{intuizione}
La regola d'oro del profiling: prima corretto, poi veloce.
Non ottimizzare prima di avere la soluzione giusta.
Non ottimizzare prima di misurare.
Il collo di bottiglia è quasi sempre dove non ti aspetti.
\end{intuizione}

\begin{esempio}[\fn{declare optimize}: dare suggerimenti al compilatore]
\begin{lstlisting}
;; Con declare optimize si spostano le priorita' del compilatore.
;; (speed 3): massima velocita'.
;; (safety 0): niente controlli di tipo a runtime.
;; (debug 0): niente informazioni di debug.
;;
;; Usare con cautela: (safety 0) puo' causare comportamenti
;; non definiti su input non validi.

(defun somma-vettore (vec)
  (declare (optimize (speed 3) (safety 0))
           (type (simple-vector) vec))
  (loop for x across vec sum x))
\end{lstlisting}
\end{esempio}

\begin{nota}
Le dichiarazioni di tipo in \fn{declare} aiutano il compilatore
a generare codice più efficiente, non solo a validare.
Su Allegro CL e SBCL, aggiungere \fn{(type fixnum n)} a una
funzione aritmetica intera può ridurre l'uso di bignum boxing
e accelerare il codice di un ordine di grandezza.
\end{nota}

% ================================================================
\section{Il progetto finale: mettere tutto insieme}
% ================================================================

Sei arrivato alla fine del percorso.
Vale la pena fare il punto sull'arco che hai percorso.

Al Capitolo~1 valutavi \fn{(+ 1 2)} nel REPL e capivi
perché le parentesi stanno dove stanno.
All'esercizio~148 hai costruito un mini-REPL:
legge espressioni dall'utente, le valuta con il tuo
evaluator, stampa il risultato.

Quella distanza è Lisp.

\argomento{Cosa ha reso possibile il mini-REPL}

Il mini-REPL del Capitolo~8 non è un esercizio isolato:
usa ogni strato del percorso:

\begin{itemize}
  \item \textbf{Capitolo~1}: il modello mentale che rende
        ovvio il codice da esaminare come dati.
  \item \textbf{Capitolo~2}: le funzioni come blocchi
        di costruzione (\fn{read-from-string}, \fn{format}).
  \item \textbf{Capitolo~3}: la ricorsione sulla struttura
        delle S-expression.
  \item \textbf{Capitolo~4}: le funzioni come valori
        nell'ambiente di valutazione.
  \item \textbf{Capitolo~5}: la alist come ambiente,
        la hash table come cache.
  \item \textbf{Capitolo~6}: la comprensione di come
        le macro trasformano il codice prima della valutazione.
  \item \textbf{Capitolo~7}: il backtracking come template
        per la ricerca.
  \item \textbf{Capitolo~8}: stream, tokenizzazione,
        il parser, l'evaluator.
\end{itemize}

\argomento{Da qui in avanti}

Un mini-REPL con le funzioni aritmetiche di base è un punto
di partenza, non un punto di arrivo.
Le direzioni naturali per estenderlo:

\begin{itemize}
  \item Aggiungere \fn{define} per creare funzioni
        nell'ambiente a runtime.
  \item Aggiungere \fn{cond} e \fn{let} (già abbozzato
        negli esercizi~145--146).
  \item Aggiungere \fn{lambda} con chiusura sull'ambiente
        (esercizio~147).
  \item Separare il codice in package, aggiungere un file \fn{.asd},
        scrivere test per ogni funzione dell'evaluator.
  \item Aggiungere la gestione delle condizioni per
        segnalare errori leggibili invece di crash.
\end{itemize}

Questo progetto, sviluppato con i principi di questo capitolo
(funzioni piccole, package, test, condizioni), è un'ottima
palestra per tutto quello che hai imparato.

\begin{workbook}
\textbf{Usare gli esercizi già fatti come banco di prova:}\\[4pt]
Non ci sono nuovi esercizi per questo capitolo.
Invece, riprendi il codice degli esercizi 121--148
e applicaci i principi di questo capitolo:

\begin{itemize}
  \item Aggiungi le docstring alle funzioni più importanti.
  \item Organizza il mini-evaluator (es. 145--148) in un package.
  \item Scrivi una suite di test \fn{assert} per ogni funzione
        del Capitolo~7 che hai implementato.
  \item Misura con \fn{time} la differenza tra \fn{fib-lento}
        e \fn{fib-veloce}: è il modo più diretto per
        capire cosa fa il profiling.
\end{itemize}
\end{workbook}

% ================================================================
\section*{Riepilogo del capitolo}
% ================================================================

\begin{tcolorbox}[
  enhanced,
  colback=csBlue!5, colframe=csBlue!40,
  arc=4pt, boxrule=0.6pt,
  left=14pt, right=14pt, top=12pt, bottom=12pt,
  before skip=16pt, after skip=16pt
]
\textbf{Quello che hai imparato:}

\medskip
\begin{itemize}
  \item \textbf{Funzioni piccole}: una funzione, una responsabilità.
        Il corpo descrivibile in una riga. Separazione dei livelli.
  \item \textbf{API}: nomi descrittivi, argomenti nell'ordine naturale,
        valore di ritorno coerente, side effect dichiarati.
  \item \textbf{Convenzioni di naming}: \fn{predicato-p}, \fn{!},
        \fn{*variabile*}, \fn{+costante+}, kebab-case.
  \item \textbf{Docstring}: prima stringa nel corpo di \fn{defun}.
        Recuperabile con \fn{documentation} e \fn{describe}.
  \item \textbf{Package}: spazio di nomi. \fn{defpackage},
        \fn{in-package}, \fn{export}. I simboli non esportati
        sono privati.
  \item \textbf{ASDF}: descrive quali file compilare e in che ordine.
        Il file \fn{.asd} è il punto di entrata del sistema.
  \item \textbf{Testing}: \fn{assert} per casi semplici,
        \fn{fiveam}/\fn{parachute} per suite strutturate.
        Scrivere il test prima chiarisce il problema.
  \item \textbf{Condizioni}: \fn{error}, \fn{warn},
        \fn{handler-case}, \fn{ignore-errors}.
        Più flessibili delle eccezioni Java grazie ai restart.
  \item \textbf{Profiling}: \fn{time} per misurare. Prima corretto,
        poi veloce. Il collo di bottiglia è quasi sempre
        dove non ti aspetti.
\end{itemize}

\medskip
\textbf{Fine del percorso.}
Hai attraversato il linguaggio dal modello mentale alla
programmazione professionale.
Il resto è pratica: scrivere codice, leggerlo, migliorarlo.
Il REPL ti aspetta.
\end{tcolorbox}
